% GENERATED FILE. Edit canonical lessons, then run npm run generate:books.
% canonical-source-hash: fnv1a64:03000f149f383767
% canonical-lessons: MR-C279-tyanantar, MR-C279-karun, MR-C279-rtu, MR-W279-letter-vocalic-r, MR-C279-pavsalyatla-divas, MR-R279-a-story-in-order

\chapter{Telling It in Order}
\label{ch:mr-telling-it-in-order}
\hlchaptermodality{279}

\begin{chapteropening}
\textbf{By the end of this chapter:} \emph{I can tell a story in order with ādhī, mag, tyānantar and śevṭī, join actions with -ūn, and read and write \mr{ऋ} in ṛtū.}

Complete MR-R279-a-story-in-order: tell a short story in order, after reading Arun's monsoon day.
\end{chapteropening}

\section[\texorpdfstring{\mr{त्यानंतर} (tyānantar)}{त्यानंतर (tyānantar)}]{\mr{त्यानंतर} (tyānantar) — first, then, after that, in the end}

\label{lesson:MR-C279-tyanantar}



\begin{quote}
Say \textquotedblleft{}then\textquotedblright{} and \textquotedblleft{}after\textquotedblright{}. (\emph{mag}, \emph{nantar}.) Then \textquotedblleft{}I used to go\textquotedblright{}. (\emph{mī jāyco}.)
\end{quote}

\subsection*{You'll want to know: first, then, after that}
\begin{quote}
\textbf{\mr{आधी} \mr{मी} \mr{चहा} \mr{घेतला}. \mr{मग} \mr{मी} \mr{पुस्तक} \mr{वाचलं}.}
\end{quote}

\begin{quote}
\emph{ādhī mī cahā ghetlā. mag mī pustak vāclaṁ.} — \textquotedblleft{}First I had tea. Then I read a book.\textquotedblright{}
\end{quote}

A story keeps its events in order with four words:

\begin{itemize}
\raggedright
  \item \textbf{\mr{आधी}} — first
  \item \textbf{\mr{मग}} — then
  \item \textbf{\mr{त्यानंतर}} (\emph{tyānantar}) — after that: \textbf{\mr{त्या}}, \textquotedblleft{}that\textquotedblright{}, and \textbf{\mr{नंतर}}, \textquotedblleft{}after\textquotedblright{}
  \item \textbf{\mr{शेवटी}} — in the end
\end{itemize}

Each one opens its sentence.

\begin{grammarlens}[title={Grammar Lens: -\mr{नंतर} after a noun}]
\begin{quote}
\textbf{\mr{शाळेनंतर} \mr{मी} \mr{घरी} \mr{आलो}.}
\end{quote}

\begin{quote}
\emph{śāḷenantar mī gharī ālo.} — \textquotedblleft{}After school I came home.\textquotedblright{}
\end{quote}

\textbf{\mr{नंतर}} also stands after a noun, in its bent form: \textbf{\mr{शाळा}} $\to$ \textbf{\mr{शाळेनंतर}} (\emph{śāḷenantar}), \textquotedblleft{}after school\textquotedblright{}. \textbf{\mr{त्यानंतर}} is the same thing done to \textbf{\mr{तो}}.
\end{grammarlens}

\subsection*{Your turn}
Say these aloud:

\begin{itemize}
\raggedright
  \item \textquotedblleft{}First I had tea. Then I read a book.\textquotedblright{} — \emph{ādhī mī cahā ghetlā. mag mī pustak vāclaṁ.}
  \item \textquotedblleft{}After that, suddenly, it rained\textquotedblright{} — \emph{tyānantar acānak pāūs paḍlā}
  \item \textquotedblleft{}After school I came home\textquotedblright{} — \emph{śāḷenantar mī gharī ālo}
\end{itemize}

\begin{tcolorbox}[breakable,colback=teal!4,colframe=teal!35!black,title={Before you move on}]
\textquotedblleft{}First\textquotedblright{}? (\textbf{\mr{आधी}}.) \textquotedblleft{}After that\textquotedblright{}? (\textbf{\mr{त्यानंतर}}.) \textquotedblleft{}In the end\textquotedblright{}? (\textbf{\mr{शेवटी}}.) \textquotedblleft{}After school\textquotedblright{}? (\textbf{\mr{शाळेनंतर}}.)
\end{tcolorbox}
\section[\texorpdfstring{\mr{करून} (karūn)}{करून (karūn)}]{\mr{करून} (karūn) — having done it, and then}

\label{lesson:MR-C279-karun}



\begin{quote}
Say \textquotedblleft{}pleased to meet you\textquotedblright{}. (\emph{bheṭūn ānand jhālā}.) \textbf{\mr{भेटून}} is \textquotedblleft{}having met\textquotedblright{}: the form this lesson teaches has been in your mouth since the first meeting.
\end{quote}

\subsection*{You'll want to know: -\mr{ऊन}, 'having done'}
\begin{quote}
\textbf{\mr{मी} \mr{घरी} \mr{जाऊन} \mr{झोपलो}.}
\end{quote}

\begin{quote}
\emph{mī gharī jāūn jhoplo.} — \textquotedblleft{}I went home and slept.\textquotedblright{}
\end{quote}

Take \textbf{-\mr{णे}} off and put \textbf{-\mr{ऊन}} on: \textbf{\mr{करणे}} $\to$ \textbf{\mr{करून}} (\emph{karūn}), \textbf{\mr{बसणे}} $\to$ \textbf{\mr{बसून}} (\emph{basūn}), \textbf{\mr{पाहणे}} $\to$ \textbf{\mr{पाहून}} (\emph{pāhūn}). A stem ending in a vowel keeps it: \textbf{\mr{जाणे}} $\to$ \textbf{\mr{जाऊन}} (\emph{jāūn}), \textbf{\mr{येणे}} $\to$ \textbf{\mr{येऊन}} (\emph{yeūn}), \textbf{\mr{घेणे}} $\to$ \textbf{\mr{घेऊन}} (\emph{gheūn}).

The tense comes once, at the end: \textbf{\mr{जाऊन} \mr{झोपलो}} is \textquotedblleft{}went and slept\textquotedblright{}. \textbf{\mr{तो} \mr{चहा} \mr{घेऊन} \mr{आला}} (\emph{to cahā gheūn ālā}), \textquotedblleft{}having taken tea he came\textquotedblright{}, is how Marathi says \textquotedblleft{}he brought tea\textquotedblright{}.

\begin{grammarlens}[title={Grammar Lens: -\mr{ऊन} or -\mr{ल्यावर}}]
\begin{quote}
\textbf{\mr{घरी} \mr{पोहोचल्यावर} \mr{मी} \mr{झोपलो}.}
\end{quote}

\begin{quote}
\emph{gharī pohocalyāvar mī jhoplo.} — \textquotedblleft{}After I reached home, I slept.\textquotedblright{}
\end{quote}

\textbf{-\mr{ऊन}} joins two actions of \textbf{one} doer, one after the other. \textbf{-\mr{ल्यावर}} says \textquotedblleft{}after\textquotedblright{}, and the doer can change: \textbf{\mr{मी} \mr{घरी} \mr{पोहोचल्यावर} \mr{आई} \mr{आली}} (\emph{mī gharī pohocalyāvar āī ālī}), \textquotedblleft{}after I reached home, Mother came\textquotedblright{}.
\end{grammarlens}

\subsection*{Your turn}
Say these aloud:

\begin{itemize}
\raggedright
  \item \textquotedblleft{}I went home and slept\textquotedblright{} — \emph{mī gharī jāūn jhoplo}
  \item \textquotedblleft{}He brought tea\textquotedblright{} — \emph{to cahā gheūn ālā}
  \item \textquotedblleft{}After I reached home, Mother came\textquotedblright{} — \emph{mī gharī pohocalyāvar āī ālī}
\end{itemize}

\begin{tcolorbox}[breakable,colback=teal!4,colframe=teal!35!black,title={Before you move on}]
\textquotedblleft{}Having done\textquotedblright{}? (\textbf{\mr{करून}}.) \textquotedblleft{}Having gone\textquotedblright{}? (\textbf{\mr{जाऊन}}.) \textquotedblleft{}He brought tea\textquotedblright{}? (\textbf{\mr{तो} \mr{चहा} \mr{घेऊन} \mr{आला}}.) Which joins two doers? (\textbf{-\mr{ल्यावर}}.)
\end{tcolorbox}
\section[\texorpdfstring{\mr{ऋतू} (ṛtū)}{ऋतू (ṛtū)}]{\mr{ऋतू} (ṛtū) — a season}

\label{lesson:MR-C279-rtu}



\begin{quote}
Before the new one: say the summer, the monsoon and the winter. (\emph{unhāḷā}, \emph{pāvsāḷā}, \emph{hivāḷā}.) Each of them is one of these.
\end{quote}

\subsection*{You'll want to know: \mr{ऋतू}}
\textbf{\mr{ऋतू}} — \emph{ṛtū} — \textquotedblleft{}a season\textquotedblright{}.

\begin{quote}
\textbf{\mr{मला} \mr{पावसाळा} \mr{हा} \mr{ऋतू} \mr{खूप} \mr{आवडतो}.}
\end{quote}

\begin{quote}
\emph{malā pāvsāḷā hā ṛtū khūp āvaḍto.} — \textquotedblleft{}I like the monsoon season very much.\textquotedblright{}
\end{quote}

Its first letter is the vowel you know as the sign \textbf{\mr{ृ}} in \textbf{\mr{नृत्य}}, standing on its own. Marathi says it \emph{ru}: \textbf{\mr{ऋतू}} is said \emph{rutū}.

\begin{cousinweb}
\textbf{\mr{ऋतू}} (\emph{ṛtū}) is a tatsama: Sanskrit \emph{ṛtu}, \textquotedblleft{}a fixed time, a season\textquotedblright{} (Monier-Williams). CDIAL has no entry for it. Everyday Marathi names three seasons; the Sanskrit calendar counts six, and \textbf{\mr{वसंत}}, \textquotedblleft{}springtime\textquotedblright{}, is one of them.
\end{cousinweb}

\subsection*{Your turn}
Say these aloud:

\begin{itemize}
\raggedright
  \item \emph{ṛtū}, said \emph{rutū}
  \item \textquotedblleft{}I like the monsoon season very much\textquotedblright{} — \emph{malā pāvsāḷā hā ṛtū khūp āvaḍto}
  \item the three seasons, one after another
\end{itemize}

\begin{tcolorbox}[breakable,colback=teal!4,colframe=teal!35!black,title={Before you move on}]
What does \textbf{\mr{ऋतू}} mean? (\textquotedblleft{}A season\textquotedblright{}.) How is it said? (\emph{rutū}.)
\end{tcolorbox}
\section[\texorpdfstring{\mr{ऋ} (ṛ (Marathi ru))}{ऋ (ṛ (Marathi ru))}]{\mr{ऋ} — the vowel letter at the start of \mr{ऋतू}}

\label{lesson:MR-W279-letter-vocalic-r}



\begin{quote}
Say \textbf{\mr{ऋतू}}, \textquotedblleft{}a season\textquotedblright{}. Then find the sign \textbf{\mr{ृ}} in \textbf{\mr{नृत्य}}, \textquotedblleft{}dance\textquotedblright{}, and say its sound: \emph{ru}.

One new character comes next: the first letter of \textbf{\mr{ऋतू}}.
\end{quote}

\subsection*{Script}
\begin{quote}
\mr{ऋ}
\end{quote}

\textbf{\mr{ऋ}} is the standing form of the sign \textbf{\mr{ृ}}. The sign hangs on a consonant (\textbf{\mr{नृत्य}}, \textbf{\mr{हृदय}}); the letter starts a word (\textbf{\mr{ऋतू}}). Its careful transliteration is \emph{ṛ}; Marathi says it \emph{ru}. Keep it apart from \textbf{\mr{र}}, the consonant \emph{ra}: \textbf{\mr{ऋ}} is a vowel letter, and it stands alone.

\subsection*{Writing: \mr{ऋ} — guided copy}
\hlblockfigure{\detokenize{figures/MR-W279-letter-vocalic-r-filmstrip.pdf}}{How \mr{ऋ} is written, stroke by stroke}

Keep \textbf{\mr{ऋ}} in front of you for this first pass.

\begin{itemize}
\raggedright
  \item \textbf{1.} start at the open upper-left tip, sweep right through the short shoulder, then turn down-left through the diagonal hook
  \item \textbf{2.} lift at the middle body's left edge, loop anticlockwise round its little ball, then run down to the right and swing the lower bowl left, round the bottom, and out to the right without lifting
  \item \textbf{3.} lift and draw the right stem top-to-bottom
  \item \textbf{4.} lift and draw the top shirorekhā left-to-right
\end{itemize}

\textbf{Pen lifts: 3.} This is one sourced learner order; hands and teachers vary. Source: Saurmandal, ‘Devanagari \mr{ऋ} stroke order.svg’, panels 1–4, Wikimedia Commons, 5 August 2023.

Copy \textbf{\mr{ऋ}} twice beside the model, saying \emph{ru} as you finish each one. Then write \textbf{\mr{ऋतू}} once.

\begin{tcolorbox}[breakable,colback=teal!4,colframe=teal!35!black,title={Before you move on}]
Which letter starts \textbf{\mr{ऋतू}}? (\textbf{\mr{ऋ}}.) How does Marathi say it? (\emph{ru}.) Where does its sign \textbf{\mr{ृ}} stand? (On a consonant: \textbf{\mr{नृत्य}}.)
\end{tcolorbox}
\section[(reading)]{(reading) — a monsoon day in the village}

\label{lesson:MR-C279-pavsalyatla-divas}



\begin{quote}
Eighty-six words, every one met before or glossed after it. The writer is Arun. He says \textbf{\mr{मला} \mr{पावसाळा} \mr{हा} \mr{ऋतू} \mr{खूप} \mr{आवडतो}}, and then he tells you why.
\end{quote}

\begin{tcolorbox}[breakable,skin=enhanced,colback=orange!6,colframe=orange!45!black,boxrule=0.5pt,arc=1mm,left=6pt,right=6pt,top=4pt,bottom=4pt,fonttitle=\bfseries,title={Read this}]
\begin{quote}
\mr{गेल्या} \mr{वर्षी} \mr{पावसाळ्यात} \mr{मी} \mr{गावाला} \mr{गेलो}. \mr{मला} \mr{पावसाळा} \mr{हा} \mr{ऋतू} \mr{खूप} \mr{आवडतो}. \mr{आधी} \mr{मी} \mr{बसने} \mr{स्टेशनला} \mr{गेलो}. \mr{मग} \mr{गाडीने} \mr{गावाला} \mr{निघालो}. \mr{गाडीत} \mr{मी} \mr{पुस्तक} \mr{वाचत} \mr{होतो}. \mr{अचानक} \mr{खूप} \mr{पाऊस} \mr{आला}. \mr{डोंगर} \mr{हिरवे} \mr{होते}. \mr{संध्याकाळी} \mr{गावाला} \mr{पोहोचल्यावर} \mr{मी} \mr{नदीवर} \mr{गेलो}. \mr{नदी} \mr{भरलेली} \mr{होती}. \mr{लहान} \mr{असताना} \mr{मी} \mr{रोज} \mr{इथे} \mr{यायचो}. \mr{मग} \mr{मी} \mr{घरी} \mr{आलो}. \mr{आई} \mr{चहा} \mr{करत} \mr{होती}. \mr{आम्ही} \mr{चहा} \mr{घेऊन} \mr{अंगणात} \mr{बसलो}. \mr{पाऊस} \mr{पडत} \mr{होता}. \mr{आईने} \mr{मला} \mr{एक} \mr{गोष्ट} \mr{सांगितली}. \mr{त्यानंतर} \mr{आम्ही} \mr{भाकरी} \mr{खाल्ली}. \mr{रात्रभर} \mr{पाऊस} \mr{पडत} \mr{होता}. \mr{शेवटी} \mr{मी} \mr{लवकर} \mr{झोपलो}. \mr{तो} \mr{दिवस} \mr{खूप} \mr{चांगला} \mr{होता}.
\end{quote}

\emph{Read it:} the passage once without stopping; do not translate

\emph{Read it:} the passage again, and sort every verb: event, scene or habit
\end{tcolorbox}

\begin{culture}
New forms: \textbf{\mr{गेल्या} \mr{वर्षी}} (\emph{gelyā varṣī}, \textquotedblleft{}last year\textquotedblright{}), the \textquotedblleft{}gone\textquotedblright{} year; \textbf{\mr{घरी}} (\emph{gharī}, \textquotedblleft{}home\textquotedblright{}); \textbf{\mr{बसने}} (\emph{basne}, \textquotedblleft{}by bus\textquotedblright{}), with the \textbf{-\mr{ने}} of a tool; and \textbf{\mr{पाऊस} \mr{आला}}, \textquotedblleft{}rain came\textquotedblright{}, which is how Marathi says a downpour began.

The scene is \textbf{\mr{वाचत} \mr{होतो}}, \textbf{\mr{करत} \mr{होती}} and \textbf{\mr{पडत} \mr{होता}}; the habit is \textbf{\mr{यायचो}}. \textbf{\mr{आईने} … \mr{सांगितली}} and \textbf{\mr{आम्ही} \mr{भाकरी} \mr{खाल्ली}} follow the thing: a story, a bhakri. Every other verb is an event.
\end{culture}

\subsection*{Your turn}
Say these aloud:

\begin{itemize}
\raggedright
  \item what Arun was doing on the train when the rain came
  \item what Mother was doing when he came home
  \item which season Arun likes, with \textbf{\mr{ऋतू}}
\end{itemize}

\begin{tcolorbox}[breakable,colback=teal!4,colframe=teal!35!black,title={Before you move on}]
Why is \textbf{\mr{पडत} \mr{होता}} in the scene form? (The rain was going on while they sat.) Which word starts the last event? (\textbf{\mr{शेवटी}}.)
\end{tcolorbox}
\section[(review)]{Review — a story in order}

\label{lesson:MR-R279-a-story-in-order}



\begin{quote}
Say \textquotedblleft{}when I came home, Mother was making tea\textquotedblright{} and \textquotedblleft{}I used to go\textquotedblright{}. (\emph{mī gharī ālo, tevhā āī cahā karat hotī}, \emph{mī jāyco}.) Then say \textbf{\mr{ऋतू}} and name its first letter: \textbf{\mr{ऋ}}, said \emph{ru}.
\end{quote}

\subsection*{Your turn}
Say these aloud:

\begin{itemize}
\raggedright
  \item First I went to the station. Then I set out for the village.
  \item After that I went home and slept.
  \item He brought tea. After school I came home.
  \item I like the monsoon season very much, with \textbf{\mr{ऋतू}}
  \item Arun's monsoon day, in four sentences, with \textbf{\mr{आधी}}, \textbf{\mr{मग}}, \textbf{\mr{त्यानंतर}} and \textbf{\mr{शेवटी}}
\end{itemize}

\begin{tcolorbox}[breakable,colback=teal!4,colframe=teal!35!black,title={Before you move on}]
\textquotedblleft{}After that\textquotedblright{}? (\textbf{\mr{त्यानंतर}}.) \textquotedblleft{}Having gone\textquotedblright{}? (\textbf{\mr{जाऊन}}.) Which joins two doers? (\textbf{-\mr{ल्यावर}}.) A season? (\textbf{\mr{ऋतू}}, \emph{ṛtū}, said \emph{rutū}.)
\end{tcolorbox}
